% TOP_PROBLEM: {"id": 600014, "title": "Baker's Dozen — A totally skew disc", "category_id": 6, "category": {"id": 6, "name": "geometry", "display_name": "Geometry", "description": "Euclidean and non-Euclidean geometry, geometric structures.", "slug": "geometry", "order_index": 6, "created_at": "2026-07-31T15:26:25.671Z"}, "status": "open", "problem_number": "AMR-005-0014", "set_id": 13}
% FIRST_POSTED: 2026-10-01
% ENTRY_KIND: statement_only
% UnsolvedMath ID 600014; source catalogue code AMR-005-0014
% !TeX program = lualatex
% Definition and sources only; this file is not an LLM solution attempt.
\documentclass[11pt,a4paper]{article}
\usepackage[margin=25mm]{geometry}
\usepackage{fontspec}
\setmainfont{lmroman10-regular.otf}[BoldFont=lmroman10-bold.otf,
 ItalicFont=lmroman10-italic.otf,BoldItalicFont=lmroman10-bolditalic.otf]
\usepackage{amsmath,amssymb,mathtools}
\usepackage{unicode-math}
\setmathfont{latinmodern-math.otf}
\AtBeginDocument{\let\setminus\smallsetminus}
\usepackage{xurl,hyperref}
\hypersetup{colorlinks=true,urlcolor=blue}
\setcounter{secnumdepth}{0}
\setlength{\parindent}{0pt}
\setlength{\parskip}{5pt}
\setlength{\emergencystretch}{3em}
\begin{document}
\section*{Baker's Dozen — A totally skew disc}\label{600014}
{\small UnsolvedMath ID 600014 · AMR-005-0014 · Geometry · AMR Open Problem Lists}
\subsection{Definitions and mathematical statement}
Let $D^3=\{u\in\mathbb R^3:\|u\|\le1\}$. A smooth embedding
$f:D^3\to\mathbb R^7$ is understood to be smooth on a neighborhood of the
disk, injective on it, and of derivative rank three. At $u$ its affine tangent
plane is
$A_u=f(u)+df_u(\mathbb R^3)$.
The image is called totally skew when, for every pair $u\ne v$,
\[
 \dim\operatorname{aff}(A_u\cup A_v)=7.
\]
Equivalently, the $7\times7$ matrix formed by concatenating the two
$7\times3$ derivative matrices and the displacement column,
$[\,df_u\mid df_v\mid f(v)-f(u)\,]$, is nonsingular. This condition includes both the relative
directions of the tangent spaces and the displacement between their base
points; disjointness of the affine planes by itself is insufficient.

The question is whether such an embedding exists. Dimension seven is the
smallest ambient dimension in which two affine three-planes can have an
affine span of dimension $2\cdot3+1$. A construction must satisfy the
condition for every distinct pair, including arbitrarily close points.
A negative answer must obstruct all smooth embeddings, not merely
polynomial examples of a fixed degree.
\subsection{Short English statement}
Can a three-dimensional disk sit smoothly in seven-dimensional space so
that every two tangent planes, together with their relative position, span
all seven dimensions?
\subsection{Sources}
\begin{itemize}
\item[S1] ULAM.AI and the UnsolvedMath Contributors, supplied problems.json, record 600014 (AMR-005-0014), AMR Open Problem Lists. Catalogue identity and source transcription; CC BY 4.0.
\url{https://huggingface.co/datasets/ulamai/UnsolvedMath}.
\item[S2] Tabachnikov - A Baker's Dozen of Problems (2015), Problem 6. Original source indexed by AMR-005-0014.
\url{https://amj.math.stonybrook.edu/html-articles/Files-2015-2024/14-01/}.
\end{itemize}
\end{document}
